\clearpage\begingroup\SourceLeaf{34b}\HistoricalHead{29}{R}{GAUSS AN SCHUMACHER.}
\begin{center}
[6.]\\[.7em]
\textsc{Gauss} an \textsc{Schumacher}. Göttingen, 12. Julius 1836.
\par\medskip\rule{.78\textwidth}{.35pt}\par
{\small [Briefwechsel zwischen \textsc{Gauss} und \textsc{Schumacher} III, Altona 1860, S. 87.]}\par
\rule{.78\textwidth}{.35pt}
\end{center}
\noindent .\ .\ .\ .\ .\ .\ .\ .\ .\ .\ .\ .\ .\ .\ .\ .\ .\ .\ .\ .\ .
\par\noindent\hspace*{\parindent}Hätte ich übrigens eine neue Waage machen zu lassen, so würde ich
wie eine Verbesserung betrachten, wenn die Tragstücke (woran die Schalen,
oder genauer die Haken mit den Schalen hängen) in ihrem unteren Oehr
eine ziemlich scharfe (dem Waagebalken parallele) \textls[120]{Schneide} hätten, während
sie jetzt abgerundet sind. Dies würde bewirken, dass die beiden äussern
Schneiden immer an denselben Stellen trügen. So wie die Sache jetzt ist,
wird bei einiger Drehung der Schalen mit den Haken, zuweilen eine etwas
mehr nach vorne, zuweilen eine etwas mehr nach hinten liegende Stelle als
wahrer Tragepunkt zu betrachten sein, was insofern relevant ist, als man
nicht darauf rechnen kann, dass die drei Schneiden \textls[120]{genau} parallel und senkrecht
gegen den Waagebalken sind.\quad .\ .\ .\ .\ .\ .\ .\ .\ .\ .\ .\ .\ .\ .\ .\ .\ .\ .\ .\ .\ .\ .\ .\ .
\clearpage\endgroup
